% GENERATED FILE. Edit canonical lessons, then run npm run generate:books.
% canonical-source-hash: fnv1a64:4bfd8e4ca9e3e22c
% canonical-lessons: KA-C501-saumya, KA-C501-mugdha, KA-C501-krura, KA-C501-nisthura, KA-C501-anakisu

\chapter{Gentle and Cruel}
\label{ch:ka-gentle-and-cruel}
\hlchaptermodality{501}

\begin{chapteropening}
\textbf{By the end of this chapter:} \emph{I can say gentle, innocent, cruel, harshly frank and to mimic mockingly.}

Complete the last lesson of chapter 501: I can say gentle, innocent, cruel, harshly frank and to mimic mockingly.
\end{chapteropening}

\section[\texorpdfstring{\kn{ಸೌಮ್ಯ} (saumya)}{ಸೌಮ್ಯ (saumya)}]{\kn{ಸೌಮ್ಯ} (saumya) — gentle; mild}

\label{lesson:KA-C501-saumya}



\begin{quote}
Before the new one: say the Kannada for attitude, then the Kannada for to imitate.
\end{quote}

\subsection*{You'll want to know: \kn{ಸೌಮ್ಯ}}
\textbf{\kn{ಸೌಮ್ಯ}} — \emph{saumya} — \textquotedblleft{}gentle; mild\textquotedblright{}.

\begin{quote}
\kn{ನಮ್ಮ} \kn{ವೈದ್ಯರು} \kn{ತುಂಬಾ} \kn{ಸೌಮ್ಯ} \kn{ಸ್ವಭಾವದವರು}.
\end{quote}

\begin{quote}
\emph{namma vaidyaru tumbā saumya svabhāvadavaru.}
\end{quote}

\begin{quote}
Our doctor has a very gentle nature.
\end{quote}

Also mild weather or a mild taste.

\begin{cousinweb}
\textbf{\kn{ಸೌಮ್ಯ}} (\emph{saumya}) is borrowed from Sanskrit \emph{saumya}, \textquotedblleft{}relating to soma\textquotedblright{}, and later \textquotedblleft{}cool and mild\textquotedblright{} (CDIAL 13620).
\end{cousinweb}

\begin{grammarlens}[title={What you've built}]
One more word for describing what you lived and felt.
\end{grammarlens}

\subsection*{Your turn}
\begin{itemize}
\raggedright
  \item \emph{Say it:} \emph{saumya}
  \item \emph{Say it:} \emph{saumya}, once more
  \item \emph{Say it:} \emph{namma vaidyaru tumbā saumya svabhāvadavaru.} — Our doctor has a very gentle nature.
  \item \emph{Recall:} say the Kannada for attitude, then the Kannada for to imitate, then say \emph{saumya} again
\end{itemize}

\begin{tcolorbox}[breakable,colback=teal!4,colframe=teal!35!black,title={Before you move on}]
What does \kn{ಸೌಮ್ಯ} mean? (\textquotedblleft{}Gentle; mild\textquotedblright{}.) Say it once more.
\end{tcolorbox}
\section[\texorpdfstring{\kn{ಮುಗ್ಧ} (mugdha)}{ಮುಗ್ಧ (mugdha)}]{\kn{ಮುಗ್ಧ} (mugdha) — innocent; naive}

\label{lesson:KA-C501-mugdha}



\begin{quote}
Before the new one: say the Kannada for to imitate, then the Kannada for gentle.
\end{quote}

\subsection*{You'll want to know: \kn{ಮುಗ್ಧ}}
\textbf{\kn{ಮುಗ್ಧ}} — \emph{mugdha} — \textquotedblleft{}innocent; naive\textquotedblright{}.

\begin{quote}
\kn{ಮಗುವಿನ} \kn{ಮುಗ್ಧ} \kn{ಪ್ರಶ್ನೆಗೆ} \kn{ಎಲ್ಲರೂ} \kn{ನಕ್ಕರು}.
\end{quote}

\begin{quote}
\emph{maguvina mugdha praśnege ellarū nakkaru.}
\end{quote}

\begin{quote}
Everyone laughed at the child's innocent question.
\end{quote}

\kn{ಮುಗ್ಧತೆ}, innocence.

\begin{cousinweb}
\textbf{\kn{ಮುಗ್ಧ}} (\emph{mugdha}) is borrowed from Sanskrit \emph{mugdha}, \textquotedblleft{}gone astray, bewildered, foolish\textquotedblright{} (CDIAL 10175). The sense softened in Kannada to \textquotedblleft{}artless, innocent\textquotedblright{}.
\end{cousinweb}

\begin{grammarlens}[title={What you've built}]
One more word for describing what you lived and felt.
\end{grammarlens}

\subsection*{Your turn}
\begin{itemize}
\raggedright
  \item \emph{Say it:} \emph{mugdha}
  \item \emph{Say it:} \emph{mugdha}, once more
  \item \emph{Say it:} \emph{maguvina mugdha praśnege ellarū nakkaru.} — Everyone laughed at the child's innocent question.
  \item \emph{Recall:} say the Kannada for to imitate, then the Kannada for gentle, then say \emph{mugdha} again
\end{itemize}

\begin{tcolorbox}[breakable,colback=teal!4,colframe=teal!35!black,title={Before you move on}]
What does \kn{ಮುಗ್ಧ} mean? (\textquotedblleft{}Innocent; naive\textquotedblright{}.) Say it once more.
\end{tcolorbox}
\section[\texorpdfstring{\kn{ಕ್ರೂರ} (krūra)}{ಕ್ರೂರ (krūra)}]{\kn{ಕ್ರೂರ} (krūra) — cruel; fierce}

\label{lesson:KA-C501-krura}



\begin{quote}
Before the new one: say the Kannada for gentle, then the Kannada for innocent.
\end{quote}

\subsection*{You'll want to know: \kn{ಕ್ರೂರ}}
\textbf{\kn{ಕ್ರೂರ}} — \emph{krūra} — \textquotedblleft{}cruel; fierce\textquotedblright{}.

\begin{quote}
\kn{ಆ} \kn{ಕ್ರೂರ} \kn{ಮನುಷ್ಯ} \kn{ನಾಯಿಯನ್ನು} \kn{ಹೊಡೆದನು}.
\end{quote}

\begin{quote}
\emph{ā krūra manuṣya nāyiyannu hoḍedanu.}
\end{quote}

\begin{quote}
That cruel man beat the dog.
\end{quote}

\kn{ಕ್ರೂರ} \kn{ಪ್ರಾಣಿ}, a fierce animal.

\begin{cousinweb}
\textbf{\kn{ಕ್ರೂರ}} (\emph{krūra}) is borrowed whole from Sanskrit \emph{krūra}, \textquotedblleft{}bloody, cruel\textquotedblright{} (CDIAL 3602).
\end{cousinweb}

\begin{grammarlens}[title={What you've built}]
One more word for describing what you lived and felt.
\end{grammarlens}

\subsection*{Your turn}
\begin{itemize}
\raggedright
  \item \emph{Say it:} \emph{krūra}
  \item \emph{Say it:} \emph{krūra}, once more
  \item \emph{Say it:} \emph{ā krūra manuṣya nāyiyannu hoḍedanu.} — That cruel man beat the dog.
  \item \emph{Recall:} say the Kannada for gentle, then the Kannada for innocent, then say \emph{krūra} again
\end{itemize}

\begin{tcolorbox}[breakable,colback=teal!4,colframe=teal!35!black,title={Before you move on}]
What does \kn{ಕ್ರೂರ} mean? (\textquotedblleft{}Cruel; fierce\textquotedblright{}.) Say it once more.
\end{tcolorbox}
\section[\texorpdfstring{\kn{ನಿಷ್ಠುರ} (niṣṭhura)}{ನಿಷ್ಠುರ (niṣṭhura)}]{\kn{ನಿಷ್ಠುರ} (niṣṭhura) — harshly frank; unsparing}

\label{lesson:KA-C501-nisthura}



\begin{quote}
Before the new one: say the Kannada for innocent, then the Kannada for cruel.
\end{quote}

\subsection*{You'll want to know: \kn{ನಿಷ್ಠುರ}}
\textbf{\kn{ನಿಷ್ಠುರ}} — \emph{niṣṭhura} — \textquotedblleft{}harshly frank; unsparing\textquotedblright{}.

\begin{quote}
\kn{ಅವನು} \kn{ನಿಷ್ಠುರವಾಗಿ} \kn{ಸತ್ಯ} \kn{ಹೇಳಿದನು}.
\end{quote}

\begin{quote}
\emph{avanu niṣṭhuravāgi satya hēḷidanu.}
\end{quote}

\begin{quote}
He told the truth bluntly.
\end{quote}

\kn{ನಿಷ್ಠುರ} \kn{ಕಟ್ಟಿಕೊ}, to make enemies by plain speaking.

\begin{cousinweb}
\textbf{\kn{ನಿಷ್ಠುರ}} (\emph{niṣṭhura}) is borrowed from Sanskrit \emph{niṣṭhura}, \textquotedblleft{}hard, cruel\textquotedblright{} (CDIAL 7505); Kannada uses it most for hard, frank words.
\end{cousinweb}

\begin{grammarlens}[title={What you've built}]
One more word for describing what you lived and felt.
\end{grammarlens}

\subsection*{Your turn}
\begin{itemize}
\raggedright
  \item \emph{Say it:} \emph{niṣṭhura}
  \item \emph{Say it:} \emph{niṣṭhura}, once more
  \item \emph{Say it:} \emph{avanu niṣṭhuravāgi satya hēḷidanu.} — He told the truth bluntly.
  \item \emph{Recall:} say the Kannada for innocent, then the Kannada for cruel, then say \emph{niṣṭhura} again
\end{itemize}

\begin{tcolorbox}[breakable,colback=teal!4,colframe=teal!35!black,title={Before you move on}]
What does \kn{ನಿಷ್ಠುರ} mean? (\textquotedblleft{}Harshly frank; unsparing\textquotedblright{}.) Say it once more.
\end{tcolorbox}
\section[\texorpdfstring{\kn{ಅಣಕಿಸು} (aṇakisu)}{ಅಣಕಿಸು (aṇakisu)}]{\kn{ಅಣಕಿಸು} (aṇakisu) — to mimic mockingly; to mock}

\label{lesson:KA-C501-anakisu}



\begin{quote}
Before the new one: say the Kannada for cruel, then the Kannada for harshly frank.
\end{quote}

\subsection*{You'll want to know: \kn{ಅಣಕಿಸು}}
\textbf{\kn{ಅಣಕಿಸು}} — \emph{aṇakisu} — \textquotedblleft{}to mimic mockingly; to mock\textquotedblright{}.

\begin{quote}
\kn{ಹುಡುಗರು} \kn{ಮುದುಕನ} \kn{ನಡಿಗೆಯನ್ನು} \kn{ಅಣಕಿಸಿದರು}.
\end{quote}

\begin{quote}
\emph{huḍugaru mudukana naḍigeyannu aṇakisidaru.}
\end{quote}

\begin{quote}
The boys mocked the old man's walk.
\end{quote}

Past (I …): \kn{ಅಣಕಿಸಿದೆ}, \emph{aṇakiside}.

Mocking by copying someone's voice or walk; \kn{ಅಣಕ}, mockery.

\begin{cousinweb}
\textbf{\kn{ಅಣಕಿಸು}} (\emph{aṇakisu}) is native Dravidian: DEDR 111 lists Kannada \emph{aṇaka}, \textquotedblleft{}mockery, derision\textquotedblright{}, and \emph{aṇakisu}, \textquotedblleft{}to mock, deride\textquotedblright{}.
\end{cousinweb}

\begin{grammarlens}[title={What you've built}]
One more word for describing what you lived and felt.
\end{grammarlens}

\subsection*{Your turn}
\begin{itemize}
\raggedright
  \item \emph{Say it:} \emph{aṇakisu}
  \item \emph{Say it:} \emph{aṇakisu}, once more
  \item \emph{Say it:} \emph{huḍugaru mudukana naḍigeyannu aṇakisidaru.} — The boys mocked the old man's walk.
  \item \emph{Recall:} say the Kannada for cruel, then the Kannada for harshly frank, then say \emph{aṇakisu} again
\end{itemize}

\begin{tcolorbox}[breakable,colback=teal!4,colframe=teal!35!black,title={Before you move on}]
What does \kn{ಅಣಕಿಸು} mean? (\textquotedblleft{}To mimic mockingly; to mock\textquotedblright{}.) Say it once more.
\end{tcolorbox}
